\chapter{Local search for feature selection}
\label{chapter:work}
The contents of this chapter are organized as follows. First, in section \ref{section:usm}, we review the local search algorithm from \citealp{Feige:2011:MNS:2079073.2079078} for unconstrained submodular optimization (USM) for non-monotone submodular functions. We review some of the claims on optimality. In subsection \ref{subsection:usm}, we review a slightly modified version of this algorithm that \citealp{NIPS2012_4689} used for maximization of the objective function $g(S)=R^2(S)+\nu f(S)$ for feature selection in linear regression setting as discussed in section \ref{section:diversity}. The optimality guarantee of this algorithm for such an objective function is reviewed. Next, we present the linear time local search algorithm from \citealp{Buchbinder:2012:TLT:2417500.2417835} and review some of the lemmas which helps us proving the optimality guarantee in section \ref{section:lusm}. Then in subsection \ref{subsection:proposed} we present an algorithm based on this for maximization of set function $g(S)$. We prove the optimality guarantee in such approach using our proposed theorem \ref{thm:mine}. Finally, in section \ref{section:gls}, we present the full greedy local search (GLS) framework presented in \citealp{NIPS2012_4689} for using this modified version of local search in the final algorithm. We also prove the optimality guarantee for the whole feature selection framework while using the linear time local search and compare with the guarantee presented in \citealp{NIPS2012_4689}.

\section{Unconstrained submodular maximization}
\label{section:usm}
Let's assume that we have a non-monotone non-negative submodular function $f:2^X\rightarrow \mathbb{R}$. The unconstrained submodular optimization is then defined as below:
\[\hat{S}=\operatorname*{arg\,max}_{S} f(S)\]
\noindent In \citealp{Feige:2011:MNS:2079073.2079078}, this has been claimed that finding an exact local maxima is hard in general (\PLS-complete). Therefore, they have derived a local search algorithm which finds an approximate $(1+\alpha)$ local maxima. First we define the term and then we review the algorithm and discuss it's optimality guarantee and time complexity.

\begin{definition}
\label{def:localoptima}
 Given a non-negative non-monotone submodular set function $f:2^X\to \mathbb{R}$, a set $S$ is called a $(1+\alpha)$-approximate local-optima, for some $\alpha\ge 0$, if
\[(1+\alpha)f(S)\ge f(S\cup \{x\}), \forall x\in S\]
\[(1+\alpha)f(S)\ge f(S\setminus \{x\}), \forall x\notin S\]
\end{definition}

\noindent Following from this definition, the next lemma (lemma 3.3 in \citealp{Feige:2011:MNS:2079073.2079078}) turns out to be useful in proving the optimality guarantee.

\begin{lemma}
\label{lemma:feige1}
 If $f$ is submodular set-function and $S$ is a $(1+\alpha)$-approximate local-optima, then for any $T\subseteq S$ or $T\supseteq S$, then \[f(T)\le(1+n\alpha)f(S)\]
\end{lemma}

\noindent Now we review the algorithm for USM for non-monotone non-negative submodular maximization.

\begin{algorithm}
LS algorithm for USM
\begin{algorithmic}[1]
\State $S\gets \operatorname*{arg\,max}_i f(X_i)$
\State $U\gets \{X_1,\cdots,X_n\}$
\While{$\exists x\in U\setminus S$ s.t. $f(S\cup\{x\})\ge (1+\epsilon/n^2)f(S)}$
  \State $S\gets S\cup\{x\}$
\EndWhile
\While{$\exists x\in S$ s.t. $f(S\setminus\{x\})\ge (1+\epsilon/n^2)f(S)}$
  \State $S\gets S\setminus\{x\}$
\EndWhile
\end{algorithmic}
\Return $S_{\text{LS}}=\operatorname*{arg\,max}_{T:T\in\{S,U\setminus S\}} f(T)$
\end{algorithm}

\noindent It is obvious from the construction that at termination, this algorithm outputs a $S_{\text{LS}}$ which is $(1+\epsilon/n^2)$ approximate local maxima. This algorithm gives an optimality guarantee of $(\frac{1}{3}-\frac{\epsilon}{n})$ and takes $\mathcal{O}(\frac{1}{\epsilon}n^3\log n)$. These proofs are not obvious and hence we provide them from \citealp{Feige:2011:MNS:2079073.2079078} for the sake of completeness.

\begin{theorem}
\label{thm:feige1}
 If $f$ is non-negative, submodular, then for \emph{any} set $T\subseteq U$ and any $\epsilon>0$, the LS algorithm forms a set $S$, s.t. $(2+2\epsilon/n)f(S)+f(U\setminus S)\ge f(T)$.
\end{theorem}
\begin{proof}
 Let's consider a set $T\subseteq U$ and let $\alpha=\epsilon/n^2$. For the set $S\cup T$, we can write $f(S\cup T)\le(1+n\alpha)f(S)$ since $S\cup T\supseteq S$ (from lemma \ref{lemma:feige1}). Similarly, $f(S\cap T)\le(1+n\alpha)f(S)$ since $S\cap T\subseteq S$. Also, from submodularity,
 \begin{align}
 f(S\cup T)+f(U\setminus S)&\ge f(S\cup T\cup (U\setminus S))+f((S\cup T)\cap(U\setminus S))\\
       &\ge f(U)+f(T\setminus S)\\
       &\ge f(T\setminus S)
\end{align}
\noindent The last line follows from the non-negativity of $f$. Similarly, 
\begin{align}
 f(S\cap T)+f(T\setminus S)&\ge f((S\cap T)\cup (T\setminus S))+f((S\cap T\cap(U\setminus S))\\
       &\ge f(T)+f(\emptyset)\\
       &\ge f(T)
\end{align}
\noindent Therefore,
\begin{align}
 2(1+n\alpha)f(S)+f(U\setminus S)&\ge f(S\cap T)+f(S\cup T)+f(U\setminus S)\\
       &\ge f(S\cap T)+f(T\setminus S)\\
       &\ge f(T)\qedhere
\end{align}
\end{proof}

\begin{theorem}
\label{thm:feige2}
 If $f$ is non-negative, submodular, then the LS algorithm gives an optimality guarantee of $(\frac{1}{3}-\frac{\epsilon}{n})$
\end{theorem}
\begin{proof}
 From theorem \ref{thm:feige1}, we can write, for the optimal set $C\subseteq U$, $(2+2\epsilon/n)f(S)+f(U\setminus S)\ge f(C)$. Let's assume that the set returned by the algorithm is $S$ which is when $f(S)\ge f(U\setminus S)$. Therefore, for this case $(2+2\epsilon/n)f(S)+f(S)\ge (2+2\epsilon/n)f(S)+f(U\setminus S)\ge f(C)$, and therefore $f(S)\ge \frac{1}{(3+2\epsilon/n)}f(C)$ which is at least $(\frac{1}{3}-\frac{\epsilon}{n})f(C)$ for all $\epsilon\ge 0$. The proof for $U\setminus S$ follows similarly.
\end{proof}

\begin{theorem}
 The LS algorithm described has complexity $\mathcal{O}(\frac{1}{\epsilon}n^3\log n)$.
\end{theorem}
\begin{proof}
 Let's assume that $f(\{\hat{x}\})=\max_i f(\{x_i\})$, \textit{i.e.} it is the maximum of all singletons. Since the function $f$ is submodular, its submodularity ratio $\gamma\ge 1$. Therefore,
 \[\min_{L\subseteq X, S:|S|\le k, S\cap L=\emptyset}\frac{\sum_{x\in S}f(L\cup\{x\})-f(L)}{f(L\cup S)-f(L)}\ge 1\]
 Without any loss of generality, we take $L=\emptyset$ and $S=C$, \textit{i.e.} the optimal set where $k = |C|$. Thus we have 
 \[\frac{\sum_{x\in C}f(\{x\})-f(\emptyset)}{f(C)-f(\emptyset)}\ge 1\]
 which leads to $n f(\{\hat{x}\})\ge |C|f(\{\hat{x}\})\ge \sum_{x\in C}f(\{x\})\ge f(C)$. Therefore, the solution set returned by the algorithm $S$ must follow $n f(\{\hat{x}\})\ge f(S)$.
 
 Now, let's assume that the algorithm ran up to $k$ iterations and in each step the function value is increased by up to a multiplicative factor of at least $(1+\epsilon/n^2)$. Therefore, for the final set, $S$, $f(S)\ge (1+\epsilon/n^2)^k f(\{\hat{x}\})$ where $\hat{x}$ was selected in the initial arg max computation over all the singletons. Thus $n f(\{\hat{x}\})\ge (1+\epsilon/n^2)^kf(\{\hat{x}\})$. Therefore, $k=\mathcal{O}(\frac{1}{\epsilon}n^2\log n)$ and since therefore the overall complexity is $\mathcal{O}(\frac{1}{\epsilon}n^3\log n)$.
\end{proof}

\subsection{Unconstrained approximate submodular maximization}
\label{subsection:usm}
In \citealp{NIPS2012_4689}, an algorithm which is very similar to the one described in previous section was used to find an approximate maxima for the objective function $g(S)=R^2(S)+\nu f(S)$. 

\begin{algorithm}
LS algorithm for maximizing $g(S)$.
\begin{algorithmic}[1]
\State $S\gets \operatorname*{arg\,max}_i f(X_i)$
\State $U\gets \{X_1,\cdots,X_n\}$
\While{$\exists x\in U\setminus S$ s.t. $f(S\cup\{x\})\ge (1+\epsilon/n^2)f(S)}$
  \State $S\gets S\cup\{x\}$
\EndWhile
\While{$\exists x\in S$ s.t. $f(S\setminus\{x\})\ge (1+\epsilon/n^2)f(S)}$
  \State $S\gets S\setminus\{x\}$
\EndWhile
\end{algorithmic}
\Return $S_{\text{LS}}=\operatorname*{arg\,max}_{T:T\in\{S,U\setminus S,U\}} g(T)$
\end{algorithm}

\noindent It is obvious that this algorithm has the same asymptotic run-time as of the previous LS algorithm for USM. We now review the optimality guarantee in this by the following theorem in \citealp{NIPS2012_4689} (theorem 7). Again, for the sake of completeness we present the proof of this theorem.

\begin{theorem}
\label{thm:usmgdas}
 The above LS algorithm gives a $\frac{1}{4+\frac{4\epsilon}{n}}$ multiplicative approximation guarantee for $\operatorname*{arg\,max}_{S} g(S)$.
\end{theorem}
\begin{proof}
 Let's assume that the optimal set for $g(S)$ is $C$. When the algorithm terminates, we therefore have a set $S$ obtained by the LS algorithm above. Now, from theorem \ref{thm:feige1}, we can write that $(2+2\epsilon/n)f(S)+f(U\setminus S)\ge f(C)$. Since $g(S)=R^2+\nu f(S)$ and $R^2$ is non-negative and $\nu\ge 0$, we have $\nu (2+2\epsilon/n)f(S)+(2+2\epsilon/n)R^2+\nu f(U\setminus S)+R^2(U\setminus S)\ge \nu f(C)$ or $(2+2\epsilon/n)g(S)+g(U\setminus S)\ge \nu f(C)$. Now, since $R^2$ is monotone as well, we have $R^2(U)\ge R^2(C)$ since $C\subseteq U$. Thus, $(2+2\epsilon/n)g(S)+g(U\setminus S)+R^2(U)\ge \nu f(C)+R^2(C)=g(C)$. Finally, we have $f$ which is non-negative, and therefore, $(2+2\epsilon/n)g(S)+g(U\setminus S)+R^2(U)+\nu f(U)\ge g(C)$ or 
 \[(2+2\epsilon/n)g(S)+g(U\setminus S)+g(U)\ge g(C)\]
 Then, let's assume that $S$ was selected in the final arg max computation. Therefore $g(S)\ge g(U\setminus S)$ and $g(S)\ge g(U)$. So, $(2+2\epsilon/n)g(S)+g(S)+g(S)\ge (2+2\epsilon/n)g(S)+g(U\setminus S)+g(U)\ge g(C)$ or $g(S)\ge \frac{1}{(4+2\epsilon/n)}g(C)\ge \frac{1}{(4+4\epsilon/n)}g(C)$. Similar argument follows for $U\setminus S$ and $U$.
\end{proof}

\section{Linear time unconstrained submodular maximization}
\label{section:lusm}
In \citealp{Buchbinder:2012:TLT:2417500.2417835}, a linear time deterministic algorithm was proposed providing an approximation guarantee $\frac{1}{3}.f(OPT)$ for USM where the objective is non-negative non-monotone. First, let us discuss briefly how the linear-time solution works.

\begin{algorithm}
Linear-time LS algorithm for maximizing $f(S)$
\begin{algorithmic}[1]
\State $X_0\gets \emptyset$
\State $Y_0\gets \{X_1,\cdots,X_n\}$
\For{$i\gets 1$ to $i\le|U|$}
  \State $a_i\gets f(X_{i-1}\cup\{x_i\})-f(X_{i-1})$
  \State $b_i\gets f(Y_{i-1}\setminus\{x_i\})-f(Y_{i-1})$
  \If{$a_i\ge b_i$}
    \State $X_i\gets X_{i-1}\cup\{x_i\}$
    \State $Y_i\gets Y_{i-1}$
  \Else
    \State $X_i\gets X_{i-1}$
    \State $Y_i\gets Y_{i-1}\setminus \{x_i\}$
  \EndIf
\EndFor
\end{algorithmic}
\Return $X_n$ or $Y_n$ (Both are same)
\end{algorithm}

\noindent The approximate optimality guarantee proof follows from two lemmas presented in \citealp{Buchbinder:2012:TLT:2417500.2417835}:
\begin{lemma}
\label{lem:ab}
 For every $i\in[1,n]$, $a_i+b_i\ge 0$.
\end{lemma}
\begin{lemma}
 For every $i\in[1,n]$, \[f(OPT_{i-1})-f(OPT_i)\le f(X_i)-f(X_{i-1})+f(Y_i)-f(Y_{i-1})\] where $OPT_i$ is defined as $(OPT\cup X_i)\cap Y_i$. Thus, $OPT_0=OPT$ and $OPT_n=X_n=Y_n$.
\end{lemma}
\noindent The final proof then works as summing up $\forall i\in[1,n]$, the inequality from this lemma.
\begin{align}
 f(OPT_0)-f(OPT_n)&\le f(X_n)-f(X_0)+f(Y_n)-f(Y_0)\\
       &\ge f(X_n)+f(Y_n)
\end{align}
\noindent and therefore, $f(S)=f(X_n)=f(Y_n)=f(OPT)/3$.

However the relation between the solution set returned by the algorithm $S$ while solving $f$, with any other set $T\subseteq U$ is not obvious from the construction of this algorithm. But, to use a linear time local search for finding $\operatorname*{arg\,max}_{S} g(S)$, we need to establish such a relation since optimal set of $f$ and $g$ are different. Therefore, we propose the following theorem establishing the base of our proofs and motivation towards the proposed algorithm.
\begin{theorem}
\label{thm:mine}
 The solution set returned by the USM linear time local search algorithm, $S$, is a $(1+\alpha)$-approximate local optima where $\alpha=\max_i(\frac{1}{f(S)}(a_i+b_i))\in [0,2]$ where $a_i = f(X_{i-1}\cup\{x_i\})-f(X_{i-1})$ and $b_i = f(Y_{i-1}\setminus\{x_i\})-f(Y_{i-1}))$
\end{theorem}
\begin{proof}
First, we notice that for any $i\in[1,n]$, the final set selected by the algorithm $S=X_n=Y_n$ would satisfy $X_i\subseteq S\subseteq Y_i$. 

For $T\subseteq S$, we can think of a chain $T=T_1\subseteq T_2\cdots \subseteq T_k=S$, where $T_j\setminus T_{j-1}=\{x_j\}$. Now for each such element $x_j\notin T_{j-1}$ and $x_j\in S$, we can write, from the property of diminishing returns of the submodular function $f$ (proposition \ref{def:second}), 
\[
 f(T_j)-f(T_{j-1})\ge f(S)-f(S\setminus \{x_j\}) \ge f(Y_{j-1})-f(Y_{j-1}\setminus \{x_j\})=-b_j
\]
Summing up, we get, $f(S)-f(T)\ge-\sum_j b_j$ or \[f(T)\le f(S)+\sum_j b_j\]
Notice that these elements have been added in the final solution set, $S$, by the algorithm. Therefore $\forall x_j\in S$, $a_j\ge b_j$. So we have \[f(T)\le f(S)+\sum_j b_j\le f(S)+\sum_j a_j\le f(S)+\sum_j(a_j+b_j)\]
Now let's consider the simplest case where $T=S\setminus \{x_k\}$ for some $k$. So we have $f(T)=f(S\setminus \{x_k\})\le f(S)+a_k+b_k$. Since $\forall i$, $a_i+b_i\ge 0$, we can write this as 
\[f(S\setminus \{x_k\})\le(1+\alpha_k)f(S)\]
for $\alpha_k=\frac{1}{f(S)}(a_k+b_k)\ge 0$.

Similarly, for $T\supseteq S$, we can consider the chain $S=T_1\subseteq T_2\cdots \subseteq T_m=T$, where $T_j\setminus T_{j-1}=\{x_j\}$ and we find that
\[
 f(T_j)-f(T_{j-1})\le f(S\cup\{x_j\})-f(S) \le f(X_{j-1}\cup\{x_j\})-f(X_{j-1})=a_j
\]
and therefore \[f(T)\le f(S)+\sum_j a_j\]
Now $\forall x_j\notin S$, we have, from the algorithm, $a_i<b_i$. So, going by similar logic we can find, for some $T=S\cup\{x_k\}$,
\[f(S\cup \{x_k\})\le(1+\alpha_k)f(S)\]
for $\alpha_k=\frac{1}{f(S)}(a_k+b_k)\ge 0$. Choosing $\alpha=\max_i(\alpha_i)=\max_i(\frac{1}{f(S)}(a_i+b_i))\ge 0$, we can write, 
\begin{itemize}
 \item for any $x_k\in S$, $f(S\setminus \{x_k\})\le(1+\alpha)f(S)$ 
 \item and for any $x_k\notin S$, $f(S\cup \{x_k\})\le(1+\alpha)f(S)$. 
\end{itemize}

\noindent Therefore $S$ is a $(1+\alpha)$-approximate local optima according to definition \ref{def:localoptima}. Now it is left to prove the upper bound of $\alpha$. To prove that, we notice that the the algorithm always increases the function value of $f$ in each iteration. This is easy to verify, since if it didn't then at least at one step in the algorithm we'd have got the gains, $a_i$ and $b_i$ both negative but according to  lemma \ref{lem:ab}, $a_i+b_i\ge 0$. Thus $f(S)\ge f(X_i)$ and $f(S)\ge f(Y_i)$.
Therefore, $\max_i \frac{1}{f(S)}(f(X_{i-1}\cup\{x_i\})-f(X_{i-1}) + f(Y_{i-1}\setminus\{x_i\})-f(Y_{i-1}))\le 2$.
\end{proof}

\subsection{Proposed algorithm for approximate submodular maximization}
\label{subsection:proposed}
For our objective function of interest, $g(S)=R^2+\nu f(S)$, we propose the following algorithm which is very similar to the one presented in \citealp{Buchbinder:2012:TLT:2417500.2417835} for USM with non-negative, non-monotone submodular objective. First we present the algorithm and then we analyze the optimality guarantee.

\begin{algorithm}
Linear-time LS algorithm for maximizing $g(S)$
\begin{algorithmic}[1]
\State $X_0\gets \emptyset$
\State $Y_0\gets \{X_1,\cdots,X_n\}$
\For{$i\gets 1$ to $i\le|U|$}
  \State $a_i\gets f(X_{i-1}\cup\{x_i\})-f(X_{i-1})$
  \State $b_i\gets f(Y_{i-1}\setminus\{x_i\})-f(Y_{i-1})$
  \If{$a_i\ge b_i$}
    \State $X_i\gets X_{i-1}\cup\{x_i\}$
    \State $Y_i\gets Y_{i-1}$
  \Else
    \State $X_i\gets X_{i-1}$
    \State $Y_i\gets Y_{i-1}\setminus \{x_i\}$
  \EndIf
\EndFor
\end{algorithmic}
\Return $S_{\text{LS}}=\operatorname*{arg\,max}_{T:T\in\{X_n,U\setminus X_n,U\}} g(T)$
\end{algorithm}

\begin{theorem}
\label{thm:usmgme}
 The proposed LS algorithm gives a $\frac{1}{4+2 n\alpha}\in[\frac{1}{4},\frac{1}{4+4n}]$ multiplicative approximation guarantee for $\operatorname*{arg\,max}_{S} g(S)$ where $\alpha=\max_i(\frac{1}{f(S)}(a_i+b_i))\in [0,2]$ where $a_i = f(X_{i-1}\cup\{x_i\})-f(X_{i-1})$ and $b_i = f(Y_{i-1}\setminus\{x_i\})-f(Y_{i-1}))$ in the above algorithm.
\end{theorem}
\begin{proof}
 First, we notice that using lemma \ref{lemma:feige1} and theorem \ref{thm:mine}, we can write, for any $T\subseteq S$ or $T\supseteq S$, then \[f(T)\le(1+n\alpha)f(S)\le (1+2n)f(S)\]
 Therefore, from lemma \ref{thm:feige1}, we can write for \emph{any} set $T\subseteq U$
 \[(2+4n)f(S)+f(U\setminus S)\ge (2+2n\alpha)f(S)+f(U\setminus S)\ge f(T)\]
 Thus, following the proof in theorem \ref{thm:usmgdas}, we can write that
 \[\max(g(S),g(U\setminus S),g(U))\ge\frac{1}{4+2 n\alpha}g(C)\ge \frac{1}{4+4n}g(C) \]
\end{proof}

\noindent In practice, the value of $\alpha$ is $\ll 2$ and also this guarantee does not depend on the $\epsilon$ parameter in the local search algorithm. In section \ref{subsection:results}, we show a comparative plot of the these two algorithms on a synthetic data set.

\section{Greedy local search algorithm for feature selection}
\label{section:gls}
In this section we review the greedy local search algorithm framework proposed by \citealp{NIPS2012_4689} where the FR algorithm (see subsection \ref{subsection:approxmax}) is used along with LS. First, we present the algorithm and review a theorem from the same proving the optimality guarantee for this algorithm. Then we propose a theorem proving the optimality guarantee if we replace the local search with the proposed local search in subsection \ref{subsection:proposed}.

\begin{algorithm}
GLS algorithm for maximizing $g(S)$
\begin{algorithmic}[1]
\State $U\gets \{X_1,\cdots,X_n\}$
\State $S_1\gets \text{FR}(U)$
\State $S'_1\gets \text{LS}(S_1)$
\State $S_2\gets \text{FR}(U\setminus S_1)$
\end{algorithmic}
\Return $S_{\text{GLS}}=\operatorname*{arg\,max}_{T:T\in\{S_1,S'_1,S_2\}} g(T)$
\end{algorithm}

\noindent First, we mention the following two results (theorem 5 and lemma 8) from \citealp{NIPS2012_4689} which helps in proving the optimality bound.

\begin{theorem}
\label{thm:fr}
 Let $T\subseteq U$ with $|T|\le k$. Then the set $S$ selected by the FR algorithm satisfies 
 \[g(S)\ge (1-e^{-\frac{\gamma_{S,2k}}{2}})g(S\cup T)\]
\end{theorem}

\begin{lemma}
\label{lem:help}
 Let's assume that we have sets $C,S_1\subseteq U$. Also, let's assume that $C'=C\setminus S_1$ and $S_2\subseteq U\setminus S_1$. Then $g(S_1\cup C)+g(S_2\cup C')+g(S_1\cap C)\ge g(C)$.
\end{lemma}

\noindent Finally we reproduce the theorem from \citealp{NIPS2012_4689} showing the optimality guarantee which works for both the local search algorithms discussed above.

\begin{theorem}
 The GLS algorithm outputs a solution set $S$ with optimality guarantee
 \[\frac{1-e^{-\frac{\gamma_{S,2k}}{2}}}{2+(1-e^{-\frac{\gamma_{S,2k}}{2}})(4+2n\alpha)}\ge \frac{1-e^{-\frac{\gamma_{S,2k}}{2}}}{6+2n\alpha}\]
\end{theorem}

\begin{proof}
 Let $C\subseteq U$, $|C|\le k$ is the optimal solution for $\operatorname*{arg\,max}_{S:|S|\le k} g(S) = R^2_Z(S) + \nu f(S)$. Applying theorem \ref{thm:fr}, we obtain that $g(S_1)\ge \lambda g(S_1\cup C)$, where $\lambda=(1-e^{-\frac{\gamma_{S,2k}}{2}})$. Now we consider two cases.\\
 Case 1: $g(S_1\cap C)\ge \delta g(C)$ for some $0\le \delta\le 1$. In this case, the local search algorithm selects a set $S'_1$ with optimality guarantee $f(S'_1)\ge\frac{\delta}{4+2n\alpha}g(C)$ following theorem \ref{thm:usmgdas} and \ref{thm:usmgme}.\\
 Case 2: $g(S_1\cap C)< \delta g(C)$. Therefore, $\lambda g(S_1\cap C) -\lambda \delta g(C)<0$. Therefore, adding this to the RHS of $g(S_1)\ge \lambda g(S_1\cup C)$, we get $g(S_1)\ge \lambda g(S_1\cup C) + \lambda g(S_1\cap C) -\lambda \delta g(C)$. Also, we have, applying \ref{thm:fr} for the second FR, $g(S_2)\ge \lambda g(S_2\cup (C\setminus S_1))$. Adding these two equations, we get $g(S_1)+g(S_2)\ge \lambda g(S_1\cup C) + \lambda g(S_1\cap C) + \lambda g(S_2\cup (C\setminus S_1)) -\lambda \delta g(C)$. Therefore, applying lemma \ref{lem:help}, we get $g(S_1)+g(S_2)\ge \lambda g(C)-\lambda \delta g(C)=\lambda(1-\delta)g(C)$. This shows that $\max(g(S_1),g(S_2))\ge \frac{\lambda(1-\delta)}{2}g(C)$.\\
 Now combining the two cases, we have an optimality guarantee
 \[\max(g(S_1),g(S'_1),g(S_2))=\max\{\frac{\delta}{4+2n\alpha},\frac{\lambda(1-\delta)}{2}\}g(C)\]
 Setting $\delta=\frac{\lambda(4+2n\alpha)}{2+\lambda(4+2n\alpha)}$ we achieve the guarantee of $\frac{1-e^{-\frac{\gamma_{S,2k}}{2}}}{2+(1-e^{-\frac{\gamma_{S,2k}}{2}})(4+2n\alpha)}$ for both these cases.
 which is at least $\frac{1-e^{-\frac{\gamma_{S,2k}}{2}}}{6+2n\alpha}$ since $1-e^{-\frac{\gamma_{S,2k}}{2}}\le 1$.
\end{proof}

\noindent Setting $\alpha=1/2n$, we obtain the optimality guarantee promised in \citealp{NIPS2012_4689} when using the local search version based on \citealp{Feige:2011:MNS:2079073.2079078}. When using linear time local search based on \citealp{Buchbinder:2012:TLT:2417500.2417835}, we therefore obtain a guarantee $\left[\frac{1-e^{-\frac{\gamma_{S,2k}}{2}}}{6},\frac{1-e^{-\frac{\gamma_{S,2k}}{2}}}{6+4n}\right]$.